\begin{exercise}[The translation zero mode of an instanton]{I-CH10-002}
For the Euclidean action
\[
  S_E[x]=\int_{-\infty}^{\infty}\dd t
  \left[\frac12\dot x^{\,2}+V(x)\right],
\]
let $X(t)$ be a finite-action instanton satisfying
$-\ddot X+V'(X)=0$.  Time-translation invariance produces the family
$X_{t_1}(t)=X(t+t_1)$.  Differentiate this family and the classical equation
to prove that $\dot X$ lies in the kernel of
\[
  \mathcal M=-\frac{\dd^2}{\dd t^2}+V''(X(t)).
\]
Also state the boundary behavior that makes this zero mode normalizable and
explain the same conclusion as a tangent-vector statement for the line of
saddles.
\end{exercise}
